\section{Proofs for the vanishing-gain instance}
\label{app:vanishing}
All expectations and almost-sure statements in this appendix, except where
marked $Q$, are under the iid Gaussian alternative $P$. Training is included
in the filtration at evaluated time zero. It has the same distribution under
$P$ and the null constructed below, so the initial density ratio is one.

\subsection{Integrated estimation error and useful moment bounds}
\label{app:vanishing-normalization}
For centered iid observations with finite fourth moment, write
$B_i=\sum_{j\le i}A_j$ and $\kappa_4=\E A_1^4-3\sigma^4$. For $i\le j$,
independent future increments give
\[
 \operatorname{Cov}(B_i^2,B_j^2)=2i^2\sigma^4+i\kappa_4.
\]
Thus $H_n=\sum_{i\le n}B_i^2/i^2$ has
$\E H_n=\sigma^2\log n+O(1)$ and $\operatorname{Var}H_n=O(\log n)$:
for the leading covariance term the inner sum over $i\le j$ contributes
$O(j)/j^2$, and the fourth-cumulant term contributes $O(\log j)/j^2$.
Chebyshev's inequality is summable along
$n_k=\lfloor\exp(k^2)\rfloor$. Borel--Cantelli gives
$H_{n_k}/\log n_k\to\sigma^2$. Positivity of the summands and
$\log n_{k+1}/\log n_k\to1$ interpolate to all $n$.
The ordinary strong laws make clipping eventually inactive and the risk
denominator converge to $m>0$. Weighted Cesaro and removal of a finite
initial prefix yield the first limit in \eqref{eq:vanishing-rate}.

With $k$ past Gaussian observations, radial/angular independence gives
\begin{align*}
 \E(\bar A_k/\widehat M_k)^2
   &=\frac1{\sigma^2(k-2)},\\
 \E|\bar A_k/\widehat M_k|^3
   &=\frac{2\sqrt2\,k^{3/2}}
       {\sqrt\pi\,\sigma^3(k-3)(k-1)(k+1)},\\
 \E(\bar A_k/\widehat M_k)^4
   &=\frac{3k}{\sigma^4(k-4)(k-2)(k+2)}.
\end{align*}
For example, rotate the Gaussian vector so its first coordinate is
$\sqrt{k}\bar A_k/\sigma$, and use its independent length and uniform
direction. These formulas apply for $k>4$. Clipping and $r\ge0$ decrease
absolute targets, hence
\begin{equation}
 \E w_t^2=O(t^{-1}),\qquad
 \sum_t\E|w_t|^3<\infty,\qquad \sum_t\E w_t^4<\infty.
 \label{eq:vanishing-moments}
\end{equation}

We repeatedly use the following martingale normalization. If a centered
adapted increment has conditional variance at most $K w_t^2$, its sum
$M_n$ satisfies $M_n/I_n\to0$ almost surely. Indeed $I_t$ is predictable and
\[
 \sum_t\frac{w_t^2}{(1+I_t)^2}
 \le\int_0^\infty\frac{dx}{(1+x)^2}=1.
\]
For constant $K$, the martingale of normalized increments is bounded in
$L^2$ and converges; the pathwise Kronecker
lemma then gives the assertion. Localization suffices when the predictable
variance bound is only almost surely summable.

\subsection{A boundary null and the attainable information rate}
Let $a_t=\gamma\sigma w_t/2$. Define $Q$ to agree with $P$ during training
and thereafter use the sequential Gaussian kernels
\begin{equation}
 A_t\mid\mathcal F_{t-1}\sim
 N\!\left(\frac{a_t\sigma}{1+a_t^2},
           \frac{\sigma^2}{(1+a_t^2)^2}\right).
 \label{eq:vanishing-null}
\end{equation}
Here $w_t$ is exactly the same measurable rule on the observed past under
both laws. The conditional raw second moment is $\sigma^2/(1+a_t^2)$,
so $\E_Q(Y_t\mid\mathcal F_{t-1})=0$. Its fourth moment is
\[
 \sigma^4\frac{3+6a_t^2+a_t^4}{(1+a_t^2)^4}\le3\sigma^4.
\]
These nondegenerate kernels give mutual absolute continuity on every finite
horizon without a small-leverage restriction. Put $R_n=dP/dQ|_{\mathcal F_n}$
and $Z_t=A_t/\sigma$. Its log increment is
\begin{align}
 \ell_t&=-\log(1+a_t^2)
  +\tfrac12\{[(1+a_t^2)Z_t-a_t]^2-Z_t^2\}\notag\\
 &=-\log(1+a_t^2)+(a_t^2+a_t^4/2)Z_t^2
          -a_t(1+a_t^2)Z_t+a_t^2/2.
 \label{eq:vanishing-lr}
\end{align}
Its conditional mean is $a_t^2/2+O(a_t^4)$ and its conditional variance is
$O(w_t^2)$, uniformly over bounded targets. Martingale normalization,
\eqref{eq:vanishing-moments} and the integrated-error limit give
$\log R_n/\log n\to\kappa$.

The linear factor in \eqref{eq:cash-linear} is $1-a_tZ_t+a_t^2Z_t^2$.
There is a finite numerical constant $K$ such that, for all real $x$,
\[
 |\log(1-x+x^2)+x-x^2/2|\le K|x|^3.
\]
Taylor expansion proves this near zero; positivity of the quadratic and its
logarithmic growth prove it away from zero. Subtracting
\eqref{eq:vanishing-lr} and expanding $\log(1+a_t^2)$ yields
\begin{equation}
 \log G_n-\log R_n
   =-\tfrac12\sum_{t\le n}a_t^2(Z_t^2-1)+\sum_{t\le n}\rho_t,
 \qquad \sum_t\E|\rho_t|<\infty.
 \label{eq:vanishing-remainder}
\end{equation}
The expectation bound follows from the Gaussian moments, bounded targets,
and \eqref{eq:vanishing-moments}. The displayed martingale is bounded in
$L^2$ since $\sum_t\E a_t^4<\infty$. Therefore the discrepancy converges
almost surely and $\sup_n\E|\log G_n-\log R_n|<\infty$, proving the second
limit in \eqref{eq:vanishing-rate}.

For the geometric Catoni assertion, $|\phi(x)-x|\le K|x|^3$ for all real
$x$, and $v_t^2-\sigma^2w_t^2=O(|w_t|^3+w_t^4)$. Substituting
$x=(\gamma/2)Y_t=-a_tZ_t+a_t^2Z_t^2$ gives, after summable remainders,
the leading log increment $-a_tZ_t+a_t^2(Z_t^2-1/2)$ of
\eqref{eq:vanishing-lr}. Its difference from $\log R_n$ thus has an
absolutely summable expected remainder, establishing the same rate.

\subsection{Upper bounds, confidence exponents and conditional benefit}
For a fixed $Q$-valid e-process $F$, any bounded stopping time $T$ satisfies
$\E_P(F_T/R_T)=\E_Q F_T\le1$. Thus $F/R$ is a $P$-valid e-process.
Its path supremum is finite almost surely by Ville's inequality, whence
$\limsup_n\log F_n/\log n\le\kappa$. The bound concerns fixed processes;
an additional argument is needed for families depending on $\alpha$.

Since $\log G_n/\log n\to\kappa>0$, first-crossing times tend to infinity
as the threshold grows, and the eventual pathwise bounds
$(\kappa-\epsilon)\log n\le\log G_n\le(\kappa+\epsilon)\log n$
sandwich their logarithms. This proves the first limit in
\eqref{eq:vanishing-cost}. The same argument along fixed-spaced review dates
proves the review-grid claim without assuming that daily crossings persist.

For an arbitrary level-$\alpha$ family $T_\alpha$, take
$N_\alpha=\lfloor\alpha^{-d}\rfloor$, $d<\kappa^{-1}$, and choose
$\eta>0$ with $\kappa d<1-\eta$. Changing measure at the bounded stopped
time gives
\begin{align*}
 P(T_\alpha\le N_\alpha)
 &\le \alpha^{\eta}
 +P\!\left\{\max_{t\le N_\alpha}\log R_t
                      >(1-\eta)\log(1/\alpha)\right\}\longrightarrow0.
\end{align*}
The maximum divided by $\log N$ tends to $\kappa$ almost surely, directly
from $\log R_n/\log n\to\kappa$; a finite initial prefix disappears in this
normalization. Finally $C_n/\log n\to\sigma^4/(2\gamma m^2)$, so composition
at $\tau_\alpha$ gives $4/\gamma$. For arbitrary families, monotonicity of
$C_n$ and the bound at $N_\alpha$ give every lower coefficient
$d\sigma^4/(2\gamma m^2)$; let $d\uparrow\kappa^{-1}$. This proves the
probability lower bound for $C$.

For realized benefits, let $M_n=J_n-C_n$. Gaussian moments give
\[
 \E_{t-1}(Y_t-\delta_t)^2
 =\sigma^2w_t^2+\frac{\gamma^2\sigma^4}{2}w_t^4\le K w_t^2,
 \qquad K=\sigma^2+\gamma^2\sigma^4L^2/2.
\]
The predictable normalization proved in Section~\ref{app:vanishing-normalization}
gives $M_n/I_n\to0$, and therefore $J_n/C_n\to1$, almost surely.
Composition at $\tau_\alpha$ proves the attained realized-benefit coefficient.
For arbitrary valid families, the preceding time lower bound implies
$T_\alpha\to\infty$ in probability. The eventual pathwise ratio then transfers
the probability lower bound from $C_{T_\alpha}$ to $J_{T_\alpha}$; it does
not require monotonicity of $J$. This completes Theorem~\ref{thm:vanishing}.

\subsection{Proof of Proposition~\ref{prop:vanishing-uniform}}
The actual conditional score second moment is
$m_{2,t}=\sigma^2w_t^2+3\gamma^2\sigma^4w_t^4/4$. The gap
\[
 d_\nu=\nu^2-\{\sigma^2L^2+3\gamma^2\sigma^4L^4/4\}
 =\gamma\sigma L^3\sqrt{M_4}
     +\gamma^2L^4(M_4-3\sigma^4)/4>0
\]
is independent of the bet. The Catoni inequality
$e^{\phi(x)}\le1+x+x^2/2$ gives
\[
 \E_{t-1}\exp\{\phi(\lambda_tY_t)-\lambda_t^2\nu^2/2\}
 \le\exp\{\lambda_t\delta_t-d_\nu\lambda_t^2/2\}
 \le\exp\{\delta_t^2/(2d_\nu)\}.
\]
Thus $\widetilde E_n^{\lambda}e^{-\mathcal B_n}$ is a $P$-supermartingale for every
predictable bet rule, where $\mathcal B_n=\sum_{t\le n}\delta_t^2/(2d_\nu)$ is common
to all such rules. The fourth-moment bound in
\eqref{eq:vanishing-moments} gives $\E\mathcal B_\infty<\infty$.
More explicitly, one may take
\[
 B_\nu=\frac3{8\gamma^2d_\nu}
 \left\{\frac13\left(\frac1{n_0-4}+\frac1{n_0-3}\right)
       +\frac1{12}\sum_{j=n_0-2}^{n_0+1}\frac1j\right\}.
\]
Indeed $\E w_t^4\le3k/[\gamma^4\sigma^4(k-4)(k-2)(k+2)]$ for
$k=n_0+t-1$; partial fractions telescope its sum to the displayed expression.
On $\{\mathcal B_\infty\le u\}$, crossing $1/\alpha$ requires the compensated
supermartingale to cross $e^{-u}/\alpha$. Ville and Markov therefore give
$P(\sup_n\widetilde E_n^{\lambda}\ge1/\alpha)
\le\alpha e^u+P(\mathcal B_\infty>u)\le\alpha e^u+B_\nu/u$.
The bound is uniform in the policy, including policies depending on $\alpha$;
take $u=\frac12\log(1/\alpha)$ to obtain its limit. No expected stopped time
is used. This specializes standard supermartingale compensation and Ville's
inequality to the constant-certificate family.

For the stated plug-in rule, the following sharper pathwise conclusion holds.
The score's conditional mean is $\gamma\sigma^2w_t^2/2$ and its conditional
variance is at most $K w_t^2$. The normalization just proved gives
$\sum_{t\le n}Y_t/I_n\to\gamma\sigma^2/2$, and hence the blockwise
past-mean bet is $\lambda_t=O(\log t/t)$ almost surely. Thus
$\sum_t\lambda_t^2<\infty$. Moreover,
\[
 \E\sum_{t\ge1}\frac{\log(t+1)}t w_t^2<\infty
\]
by \eqref{eq:vanishing-moments}. The same series is finite almost surely,
which implies $\sum_t\lambda_t\delta_t<\infty$. The elementary Catoni bound
$|\phi(x)-x|\le x^2/2$ makes the conditional transformed drift absolutely
summable, using $\E_{t-1}Y_t^2\le\nu^2$. The centered transformed increments
have conditional variance at most $\lambda_t^2\nu^2$, so their martingale
converges by localization. The penalty sum is finite as well. The log capital
therefore converges to a finite real limit. Its supremum is finite almost
surely, which proves the asserted high-confidence crossing limit.
